% !TeX root = ../../Book3.tex
% 纲要标题：### 2.4 局部—整体原则
% 纲要位置：工作记录/计划纲要.md，第 1862 行。
% 纲要细目（写作范围）：
% * 元素和模的消失判据
% * 单射、满射及同构的局部检测
% * 有限基本开覆盖与邻域有限性
% * 有限基本开覆盖上的相容局部元素黏合

\section{局部—整体原则}
\label{b3:sec:0204}

在一个素理想处看不见某个元素，只说明有一个不属于该素理想的乘子零化它；若在每个极大理想处都看不见它，这些乘子生成的理想便不能包含于任何极大理想，只能是单位理想。由此可以检测元素是否为零、是否属于给定子模，以及给定同态是否单射、满射或同构。这些判据不要求模有限生成；把点处的有限生成元延伸到开邻域，则需要另行考察有限性。

\subsection{元素与模的消失判据}
\label{b3:subsec:0204:01}

\begin{theorem}\label{b3:thm:local-zero-detection}
设 \(R\) 为交换含幺环，\(M\) 为 \(R\) 模，\(m\in M\)。下列条件等价：
\begin{enumerate}
\item \(m=0\)；
\item 对每个素理想 \(\mathfrak p\)，\(m/1=0\) 在 \(M_{\mathfrak p}\) 中成立；
\item 对每个极大理想 \(\mathfrak m\)，\(m/1=0\) 在 \(M_{\mathfrak m}\) 中成立。
\end{enumerate}
因此 \(M=0\) 当且仅当所有 \(M_{\mathfrak m}\) 为零，也当且仅当所有 \(M_{\mathfrak p}\) 为零。
\end{theorem}
\begin{proof}
\((1)\Rightarrow(2)\Rightarrow(3)\) 直接成立，其中第二次蕴含用到极大理想是素理想。为证明 \((3)\Rightarrow(1)\)，设 \(m\ne0\)，则其零化子 \(\operatorname{Ann}_R(m)\) 是真理想，故包含于某个极大理想 \(\mathfrak m\)。若 \(m/1=0\) 在 \(M_{\mathfrak m}\) 中成立，\eqref{b3:eq:localized-element-zero}给出 \(s\notin\mathfrak m\) 且 \(sm=0\)，与 \(\operatorname{Ann}_R(m)\subseteq\mathfrak m\) 矛盾。因此第三个条件推出第一个。对 \(M\) 的每个元素应用该结论，即得模的断言。若 \(R\) 为零环，幺性模只能为零，空集上的检验也与结论相符。
\end{proof}

\begin{corollary}\label{b3:cor:submodule-local-detection}
设 \(R\) 为交换含幺环，\(M\) 为 \(R\) 模，\(N\subseteq M\) 为子模，\(m\in M\)。则
\[
m\in N\quad\Longleftrightarrow\quad
m/1\in N_{\mathfrak m}\text{ 对每个极大理想 }\mathfrak m.
\]
特别地，两个子模 \(N,L\subseteq M\) 相等，当且仅当它们在每个极大理想处局部化后相等。
\end{corollary}
\begin{proof}
由\cref{b3:thm:localization-exact}，\((M/N)_{\mathfrak m}\cong M_{\mathfrak m}/N_{\mathfrak m}\)。所以条件恰好表示 \(m+N\) 在每个局部化中为零；应用\cref{b3:thm:local-zero-detection}即得 \(m+N=0\)。对子模的每个元素分别应用包含判据，得到相等判据。
\end{proof}

例如整环 \(R\) 在分式域 \(K\) 内满足
\[
R=\bigcap_{\mathfrak m\text{ 极大}}R_{\mathfrak m}.
\]
确实，对 \(z\in K\)，分母理想 \(I_z=\{a\in R:az\in R\}\) 若为真理想，选包含它的极大理想 \(\mathfrak m\)。若 \(z\in R_{\mathfrak m}\)，则 \(z=b/s\)、\(s\notin\mathfrak m\)，于是 \(s\in I_z\)，矛盾。故处处属于局部环便迫使 \(I_z=R\)、\(z\in R\)。

\subsection{单射、满射与同构的局部检测}
\label{b3:subsec:0204:02}

\begin{theorem}\label{b3:thm:local-morphism-detection}
设 \(R\) 为交换含幺环，\(M,N\) 为 \(R\) 模。对任意 \(R\) 模同态 \(f:M\rightarrow N\)，单射、满射及同构这三种性质，均可在所有极大理想处的 \(f_{\mathfrak m}\) 上检测，也可在所有素理想处检测。更一般地，若 \(L,M,N\) 为 \(R\) 模，则复形
\[
L\xrightarrow{u}M\xrightarrow{v}N,\qquad vu=0,
\]
在 \(M\) 处正合，当且仅当每个极大理想处的局部化复形在中间项正合。
\end{theorem}
\begin{proof}
正合性定理给出
\[
(\ker f)_{\mathfrak m}\cong\ker f_{\mathfrak m},\qquad
(\operatorname{coker}f)_{\mathfrak m}\cong
\operatorname{coker}f_{\mathfrak m}.
\]
用\cref{b3:thm:local-zero-detection}分别检测这两个模是否为零，便得到单射与满射的判据；两者同时成立等价于同构。对复形同样考虑
\(H=(\ker v)/\operatorname{im}u\)。由局部化与核、像及商相容，其局部化恰为局部复形的中间同调。因此所有局部复形正合等价于所有 \(H_{\mathfrak m}=0\)，也就等价于 \(H=0\)。改用所有素理想的证明完全相同。
\end{proof}

上述结论检测的是一个已经给定的映射。若只知道每个 \(M_{\mathfrak m}\) 与 \(N_{\mathfrak m}\) 各自同构，却没有一个共同的全局映射，就不能直接应用该定理。局部同构之间是否相容，是另一项需要处理的数据。

\ProseParagraphBreak
给定一组元素 \(m_1,\ldots,m_r\in M\)，它们生成 \(M\)，当且仅当其像生成每个 \(M_{\mathfrak m}\)：把判据应用于 \(R^r\rightarrow M\)、\(e_i\mapsto m_i\) 的满射性即可。这里生成元是全局预先选定的；不能偷换成“每个点都能另选有限个局部生成元”。

\subsection{有限基本开覆盖与邻域有限性}
\label{b3:subsec:0204:03}

\begin{proposition}\label{b3:prop:principal-cover-finiteness}
设 \(R\) 为交换含幺环，\(M\) 为 \(R\) 模，\(f_1,\ldots,f_r\in R\)（\(r\ge1\)）生成单位理想。
\begin{enumerate}
\item 若 \(M_{f_i}=0\) 对每个 \(i\) 成立，则 \(M=0\)。
\item 若每个 \(M_{f_i}\) 都是有限生成 \(R_{f_i}\) 模，则 \(M\) 有限生成。
\end{enumerate}
\end{proposition}
\begin{proof}
对 \(m\in M\)，第一项的假设给出整数 \(n_i\ge1\) 使 \(f_i^{n_i}m=0\)。理想 \((f_1^{n_1},\ldots,f_r^{n_r})\) 仍为单位理想：否则它包含于某个极大理想，该极大理想由素性包含每个 \(f_i\)，与原假设矛盾。将 \(1\) 写成这些幂的线性组合，便得 \(m=0\)。

对第二项，把每个局部模的有限生成元写成 \(m_{ij}/f_i^{a_{ij}}\)。取由全部有限个分子 \(m_{ij}\) 生成的子模 \(N\)。由于 \(f_i\) 已可逆，\(N_{f_i}=M_{f_i}\)；正合性给出 \((M/N)_{f_i}=0\)。第一项应用于 \(M/N\)，即得 \(M=N\)。
\end{proof}

第一章中 \(D(f_i)\) 覆盖素谱恰好等价于 \((f_1,\ldots,f_r)=R\)。所以该命题说明，在一组有限基本开覆盖上给出有限生成性，足以得到全局有限生成性。若每个点有这样的基本开邻域，还可由素谱拟紧性抽取有限子覆盖。

\ProseParagraphBreak
仅要求每个茎 \(M_{\mathfrak p}\) 有限生成则不够。例如
\[
M=\bigoplus_{p\text{ 为素数}}\mathbb Z/p\mathbb Z
\]
不是有限生成整数模，因为有限个元素只涉及有限个直和分量；但 \(M_{(p)}\cong\mathbb Z/p\mathbb Z\)，而 \(M_{(0)}=0\)。另一方面，对每个非零整数 \(f\)，
\[
M_f\cong\bigoplus_{p\nmid f}\mathbb Z/p\mathbb Z.
\]
右边仍有无限多个非零分量，因而不是有限生成 \(\mathbb Z_f\) 模。因此所有茎都有限生成，并不能保证存在局部模有限生成的基本开邻域。这与\cref{b3:thm:support-finite}对有限生成性的要求相呼应。

\ProseParagraphBreak
此外，“在局部环上检验”不能随意改成“张量到剩余域上检验”。取域 \(k\)。在 \(R=k[t]_{(t)}\) 上，商映射 \(R/(t^2)\rightarrow R/(t)\) 不是同构，但张量到唯一的剩余域 \(k\) 后成为 \(k\) 上的恒等映射。下一章的 Nakayama 引理能在有限生成条件下用剩余向量空间检测生成性；它并不使剩余域张量变成一般的正合操作。

\subsection{相容局部元素的黏合}
\label{b3:subsec:0204:04}

前面的检测定理从已经给定的全局元素或同态出发。若先在有限个基本开集上给出模元素，还须检查它们在交集上相等，才能反过来构造全局元素。这个构造仍只需局部化的分式零判据与单位理想条件。

\begin{lemma}\label{b3:lem:localization-sheaf-gluing}
设 \(A\) 为交换含幺环，\(M\) 为 \(A\) 模，\(r\ge1\)，\(f_1,\ldots,f_r\in A\)，且 \(D(f_1),\ldots,D(f_r)\) 覆盖 \(\operatorname{Spec}A\)。在下列映射中，\(M\rightarrow\prod_{i=1}^rM_{f_i}\) 将 \(m\) 送到 \((m/1)_i\)；定义 \(\delta((s_i))_{ij}\) 时，先通过自然局部化映射将 \(s_i\in M_{f_i}\) 与 \(s_j\in M_{f_j}\) 送入 \(M_{f_if_j}\)，再取两者之差。则
\[
0\longrightarrow M\longrightarrow\prod_{i=1}^rM_{f_i}
\xrightarrow{\delta}\prod_{1\le i,j\le r}M_{f_if_j},
\qquad \delta((s_i))_{ij}=s_i-s_j
\]
在 \(M\) 及 \(\prod_{i=1}^rM_{f_i}\) 处正合。换言之，在各交集上相等的一组局部元素来自唯一的 \(M\) 中元素。
\end{lemma}
\begin{proof}
若 \(m\) 在各 \(M_{f_i}\) 中为零，\cref{b3:prop:principal-cover-finiteness}的第一项便给出 \(m=0\)，所以第一个映射单射。它的像在各交集上的限制相等，因而包含于 \(\ker\delta\)。

现设 \((s_i)\in\ker\delta\)。把每个 \(s_i\) 写成 \(u_i/f_i^{n_i}\)，其中 \(u_i\in M\)、\(n_i\ge0\)。因为指标 \(i\) 只有有限多个，可以选整数 \(N\ge1\)，使 \(N\ge n_i\) 对每个 \(i\) 成立。令 \(t_i=f_i^{N-n_i}u_i\)，则
\[
t_i/1=f_i^Ns_i\quad\text{在 }M_{f_i}\text{ 中成立}.
\]
在 \(M_{f_j}\) 中考虑 \(t_i/1-f_i^Ns_j\)。它进一步在 \(f_i\) 处局部化后为零，因为 \((M_{f_j})_{f_i}\cong M_{f_if_j}\)，且 \(s_i,s_j\) 在这一交集上的像相等。由分式零判据，对每对 \((i,j)\) 存在整数 \(L_{ij}\ge0\)，使 \(f_i^{L_{ij}}(t_i/1-f_i^Ns_j)=0\)。指标对也只有有限多个，故可选 \(L\ge L_{ij}\) 对所有 \((i,j)\) 成立。于是
\[
f_i^Lt_i/1=f_i^{N+L}s_j\quad\text{在 }M_{f_j}\text{ 中成立}.
\]
令 \(N'=N+L\)、\(t_i'=f_i^Lt_i\)，便得到 \(t_i'/1=f_i^{N'}s_j\) 在每个 \(M_{f_j}\) 中同时成立。此时 \((f_1^{N'},\ldots,f_r^{N'})=A\)：若这个理想包含于某个极大理想 \(\mathfrak m\)，由 \(N'\ge1\) 和素性可知所有 \(f_i\in\mathfrak m\)，这与 \(D(f_i)\) 覆盖全谱矛盾。因此可以选 \(a_1,\ldots,a_r\in A\)，使 \(\sum_i a_if_i^{N'}=1\)。取
\[
m=\sum_{i=1}^r a_it_i'\in M.
\]
对每个 \(j\)，在 \(M_{f_j}\) 中有
\[
m/1=\sum_{i=1}^r a_i(t_i'/1)
=\left(\sum_{i=1}^r a_if_i^{N'}\right)s_j=s_j.
\]
所以 \((s_i)\) 属于第一个映射的像，唯一性则由该映射的单射性得到。
\end{proof}

第二十四章构造仿射概形的结构层时，将用这条引理把基本开集上的相容截面拼成一个截面。
